% TOP_PROBLEM: {"id": 30001256, "title": "De Giorgi Conjecture for the Allen–Cahn Equation", "category_id": 9, "category": {"id": 9, "name": "pde", "display_name": "Partial Differential Equations", "description": "PDEs and their applications in physics and geometry.", "slug": "pde", "order_index": 9, "created_at": "2026-07-31T15:26:25.671Z"}, "status": "open"}
% FIRST_POSTED: 2026-10-01
% ATTEMPT_STATUS: unresolved
% !TeX program = lualatex
% Definition and sources only; this file is not an LLM solution attempt.
\documentclass[11pt]{article}
\usepackage[margin=25mm]{geometry}
\usepackage{fontspec}
\setmainfont{Latin Modern Roman}
\usepackage{amsmath,amssymb,mathtools}
\usepackage{unicode-math}
\setmathfont{Latin Modern Math}
\usepackage{xurl,hyperref}
\hypersetup{colorlinks=true,urlcolor=blue}
\setcounter{secnumdepth}{0}
\setlength{\parindent}{0pt}
\setlength{\parskip}{5pt}
\setlength{\emergencystretch}{3em}
\begin{document}
\section*{\texorpdfstring{De Giorgi Conjecture for the Allen–Cahn Equation}{De Giorgi Conjecture for the Allen–Cahn Equation}}\label{30001256}
\subsection{Definitions and mathematical statement}
Fix $4\le N\le8$ and let $u\in C^2({\mathbb{R}}^N)$ take values in $(-1,1)$. Assume
\[
 \Delta u=u^3-u,\qquad \partial_eu>0\text{ everywhere for some unit vector }e.
\]
The assertion is that $u$ is one-dimensional: there exist a unit vector $a$ and $g:{\mathbb{R}}\to(-1,1)$ with $u(x)=g(a\cdot x)$. This describes solutions on all of Euclidean space, not a bounded domain, and strict directional monotonicity is the selected hypothesis. Limits $u\to\pm1$ in the monotone direction, energy-growth restrictions, and minimizing properties are not assumed unless independently derived. Adding them would restrict the conjecture.
\par\smallskip\textit{Formulation and concept references:} \href{https://arxiv.org/abs/2503.21245}{[S1]}.

\subsection{Short English statement}
In dimensions four through eight, must every bounded entire Allen–Cahn solution that is strictly monotone in one direction depend on only one linear coordinate?
\subsection{Research attempt: a ground-state identity and a conditional rigidity theorem}
\textbf{Status and source audit.} We derive stability from the stated strict monotonicity and prove one-dimensionality under an additional explicit integral condition. That condition is not part of the question and is not established in dimensions four through eight. The linked source [S1] studies a free-boundary Allen--Cahn/Bernoulli problem in dimension three; its title and abstract do not constitute a solution of the classical entire equation selected here.

\subsubsection{1. Monotonicity produces a positive ground state}
Let $q=\partial_eu>0$. Standard interior elliptic bootstrapping for the polynomial equation makes $u$ smooth, so differentiating is legitimate. With $V=3u^2-1$,
\[-\Delta q+Vq=0.\]
For $\phi\in C_c^\infty(\mathbb R^N)$, write $\phi=q\eta$. Expansion and integration by parts give
\[\int (|\nabla\phi|^2+V\phi^2)
=\int q^2|\nabla\eta|^2+\int\{\nabla q\cdot\nabla(q\eta^2)+Vq^2\eta^2\}
=\int q^2|\nabla(\phi/q)|^2\ge0.\tag{1}\]
There is no assumption that $q$ is bounded away from zero at infinity: division occurs only on the compact support, where positivity and continuity suffice. Equation (1) proves stability under compactly supported variations for every solution in the target.

\subsubsection{2. Quotients of derivatives and a quantitative sufficient condition}
For each coordinate let $w_i=\partial_i u$ and $\sigma_i=w_i/q$. Both $w_i$ and $q$ solve the same linearized equation, hence
\[\operatorname{div}(q^2\nabla\sigma_i)=q\Delta w_i-w_i\Delta q=0.\tag{2}\]
Testing (2) with $\eta^2\sigma_i$, where $\eta$ has compact support, yields
\[\int q^2\eta^2|\nabla\sigma_i|^2
\le4\int q^2\sigma_i^2|\nabla\eta|^2=4\int w_i^2|\nabla\eta|^2.\tag{3}\]
Suppose additionally that, for every $i$,
\[\int_{B_{2R}\setminus B_R}w_i^2=o(R^2)\qquad(R\to\infty).\tag{4}\]
Choose cutoffs equal to one on $B_R$, zero outside $B_{2R}$, with gradient bounded by $C/R$. The right side of (3) tends to zero. For each fixed ball the integral of $q^2|\nabla\sigma_i|^2$ is therefore zero, so $\sigma_i$ is constant. Write these constants as a vector $c$. Then $\nabla u=qc$ and $e\cdot c=1$, so $c\ne0$. Every derivative tangent to $c^\perp$ vanishes. Integrating along lines in that hyperplane proves $u(x)=g(a\cdot x)$ with $a=c/|c|$.

\subsubsection{3. Classification after the reduction}
The resulting $g$ satisfies $g''=g^3-g$, is bounded in $(-1,1)$, and, orienting $a$ appropriately, is strictly increasing. Its limits $\ell_-<\ell_+$ exist. The ODE and bounded monotonicity imply $g'\to0$ and $\ell_\pm^3-\ell_\pm=0$: otherwise bounded second derivative together with a persistent nonzero slope or acceleration contradicts the finite limits. Multiplication by $g'$ gives
\[\tfrac12(g')^2-\tfrac14(1-g^2)^2=C.\]
The endpoint potentials must agree. Among the possible limits $-1,0,1$, the only distinct ordered pair with equal potential is $(-1,1)$. Consequently $C=0$, $g'=(1-g^2)/\sqrt2$, and separation of variables gives
\[g(t)=\tanh((t-t_0)/\sqrt2).\]
Thus the additional hypothesis (4) gives the full expected profile, not merely a weaker local flatness statement.

\subsubsection{4. Verification and gap}
The sign in the linearized operator is fixed by $\Delta u=u^3-u$, and the cancellation in (1) uses $-\Delta q+Vq=0$. Inequality (3) uses $2ab\le\tfrac12a^2+2b^2$. Condition (4) is sufficient, not necessary: even planar fronts in higher dimensions can have derivative energy of order $R^{N-1}$ and need not satisfy it. In particular, a uniform pointwise gradient bound only gives an $O(R^N)$ volume bound and cannot imply (4). We have no argument forcing the derivative ratios to be constant under the original hypotheses alone. That missing Liouville or geometric rigidity step remains the conjectural part.

\subsection{Sources}
\begin{itemize}
\item[S1]Hardy Chan, Xavier Fernández-Real, Alessio Figalli, and Joaquim Serra, arXiv:2503.21245.
\url{https://arxiv.org/abs/2503.21245}.
\textit{Formulation source.}
\end{itemize}
\end{document}
